\documentclass[
  bibliography=totoc,     % Literatur im Inhaltsverzeichnis
  captions=tableheading,  % Tabellenüberschriften
  titlepage=firstiscover, % Titelseite ist Deckblatt
]{scrartcl}

%irgendetwas mit Tabelen und Figuren anders Nummerieren 
\usepackage{chngcntr}
\usepackage{longtable} 

\usepackage{titling}

%Textdatein einfügen
\usepackage{verbatim}

% Paket float verbessern
\usepackage{scrhack}

% Warnung, falls nochmal kompiliert werden muss
\usepackage[aux]{rerunfilecheck}

% unverzichtbare Mathe-Befehle
\usepackage{amsmath}
% viele Mathe-Symbole
\usepackage{amssymb}
% Erweiterungen für amsmath
\usepackage{mathtools}

% Fonteinstellungen
\usepackage{fontspec}
% Latin Modern Fonts werden automatisch geladen
% Alternativ zum Beispiel:
%\setromanfont{Libertinus Serif}
%\setsansfont{Libertinus Sans}
%\setmonofont{Libertinus Mono}

% Wenn man andere Schriftarten gesetzt hat,
% sollte man das Seiten-Layout neu berechnen lassen
\recalctypearea{}

% deutsche Spracheinstellungen
\usepackage[ngerman]{babel}


\usepackage[
  math-style=ISO,    % ┐
  bold-style=ISO,    % │
  sans-style=italic, % │ ISO-Standard folgen
  nabla=upright,     % │
  partial=upright,   % ┘
  warnings-off={           % ┐
    mathtools-colon,       % │ unnötige Warnungen ausschalten
    mathtools-overbracket, % │
  },                       % ┘
]{unicode-math}

% traditionelle Fonts für Mathematik
\setmathfont{Latin Modern Math}
% Alternativ zum Beispiel:
%\setmathfont{Libertinus Math}

\setmathfont{XITS Math}[range={scr, bfscr}]
\setmathfont{XITS Math}[range={cal, bfcal}, StylisticSet=1]

% Zahlen und Einheiten
\usepackage[
  locale=DE,                   % deutsche Einstellungen
  separate-uncertainty=true,   % immer Unsicherheit mit \pm
  per-mode=symbol-or-fraction, % / in inline math, fraction in display math
]{siunitx}

\DeclareSIUnit{\channel}{Channel}
\DeclareSIUnit{\year}{a}

% chemische Formeln
\usepackage[
  version=4,
  math-greek=default, % ┐ mit unicode-math zusammenarbeiten
  text-greek=default, % ┘
]{mhchem}

% richtige Anführungszeichen
\usepackage[autostyle]{csquotes}

% schöne Brüche im Text
\usepackage{xfrac}

% Standardplatzierung für Floats einstellen
\usepackage{float}
\floatplacement{figure}{htbp}
\floatplacement{table}{htbp}

% Floats innerhalb einer Section halten
\usepackage[
  section, % Floats innerhalb der Section halten
  below,   % unterhalb der Section aber auf der selben Seite ist ok
]{placeins}

% Seite drehen für breite Tabellen: landscape Umgebung
\usepackage{pdflscape}

% Captions schöner machen.
\usepackage[
  labelfont=bf,        % Tabelle x: Abbildung y: ist jetzt fett
  font=small,          % Schrift etwas kleiner als Dokument
  width=0.9\textwidth, % maximale Breite einer Caption schmaler
]{caption}
% subfigure, subtable, subref
\usepackage{subcaption}

% Grafiken können eingebunden werden
\usepackage{graphicx}
\usepackage{wrapfig}

% schöne Tabellen
\usepackage{booktabs}
\usepackage[table]{xcolor}

% Verbesserungen am Schriftbild
\usepackage{microtype}

% Literaturverzeichnis
\usepackage[
  backend=biber,
  sorting=none
]{biblatex}
% Quellendatenbank
\addbibresource{lit.bib}
\addbibresource{programme.bib}

% Hyperlinks im Dokument
\usepackage[
  german,
  unicode,        % Unicode in PDF-Attributen erlauben
  pdfusetitle,    % Titel, Autoren und Datum als PDF-Attribute
  pdfcreator={},  % ┐ PDF-Attribute säubern
  pdfproducer={}, % ┘
]{hyperref}
% erweiterte Bookmarks im PDF
\usepackage{bookmark}

% Trennung von Wörtern mit Strichen
\usepackage[shortcuts]{extdash}

%\setcounter{tocdepth}{3} % + subsubsections



\author{%
  Benedikt Lütke Lanfer \\%
  \href{mailto:benedikt.luetkelanfer@tu-dortmund.de}{benedikt.luetkelanfer@tu-dortmund.de}%
  \and%
  Enno Wellmann \\%
  \href{mailto:enno.wellmann@tu-dortmund.de}{enno.wellmann@tu-dortmund.de}%
}
\publishers{TU Dortmund – Fakultät Physik}


\newcommand*\diff{\mathop{}\!\mathrm{d}}

\NewDocumentCommand \OverfullCenter {+m} {
\noindent\makebox[\linewidth]{#1} }

\usepackage{adjustbox}


% %Tabellen und Figuren Einstellung
% \counterwithout{table}{section}
% \counterwithout{figure}{section}
% \renewcommand{\thetable}{\Roman{table}}
% \renewcommand{\thefigure}{\Roman{figure}}

%Richtiges Einrücken
\setlength{\parindent}{0pt}


\title{V18:\\ Germanium Detektor}
\author{Benedikt Lütke Lanfer \and Enno Wellmann}
\date{06. Mai 2024}
\publishers{TU Dortmund – Fakultät Physik}

\begin{document}
\begin{titlingpage}
    \begin{center}
        \begin{Huge}
            \textbf{\thetitle\\}
        \end{Huge}
    \end{center}
    \vspace{4cm}
    \includegraphics[width=\textwidth]{Bilder/Logo_TU.png} \\
    \vspace{4cm}
    \begin{center}
        \begin{huge}
            \theauthor\\
        \end{huge}
        \vspace{0.5cm}
        \begin{Large}
            benedikt.luetkelanfer@tu-dortmund.de\\
            enno.wellmann@tu-dortmund.de\\
            \vspace{1.4cm}
            Bearbeitet: \today\\
            Durchgeführt: \thedate\\
            TU Dortmund – Fakultät Physik\\
        \end{Large}
    \end{center}
\end{titlingpage}
\tableofcontents
\newpage

\section{Zielsetzung}
In diesem Versuch wird ein Germanium Energiedetektor mit einer geeichten \textbf{Eu} Quelle kalibriert.
Die

Stichworte:
\begin{itemize}
    \item Vollenergienachweiswahrscheinlichkeit
    \item Linien des Eu Spektrums 
    \item Halbwertsbreite
    \item Compton Kante
    \item 
\end{itemize}



%---------------------------------------------------------------------------------------------------------------------------------------------------------------%

\section{Theorie}
\subsection[]{Wechselwirkung von Strahlung mit Materie}
Aus \cite{book:knoll}.
Die hauptsächlichen Wechselwirkungen von Photonen mit Materie die photoelektrische Absorption,
die Compton Streuung und die paar Bildung.
Bei der photoelektrischen  Absorption wechselwirkt das Photon mit einem im Atom gebundenen Elektron und befördert es aus seinem gebundenen Zustand.
Die Energie des Photons wird danach von dem Elektron getragen und das Atom emittiert Röntgenstrahlung.
Diese Interaktion findet vor allem bei niederenergetischen $\gamma$ Photonen statt. 
Eine grobe Annäherung für das Verhalten der Wechselwirkungswahrscheinlichkeit
ist der Therm 
\begin{align}
    \tau = \text{constant} \cdot \frac{Z^n}{E_{\gamma}^{3.5}}
\end{align} %TODO herausfinden, was die Bedeutung von \tau ist.
wobei $n$ zwischen 4 und 5 liegt.

Die Compton Streuung ist oft die häufigste Interaktion der Strahlung mit dem Material.
Sie ist (abhängig von $Z$) dominant bei Energien im Bereich von etwa \qtyrange{0.5}{10}{\MeV}
Ihr Differentieller Wirkungsquerschnitt lautet
\begin{align}
    \frac{d\sigma}{d \Omega} = Z r_0^2 \left(\frac{1}{1+\alpha(1-\cos\theta)}\right)^2 %
    \left(\frac{1+ \cos^2\theta}{2}\right)%
    \left(1+ \frac{\alpha^2(1-\cos \theta)^2}{(1+\cos^2 \theta)[1+\alpha(1-\cos(\theta))]}\right)
    \label{eq:wq_compton}
\end{align}% Diese Formel ist zu kompliziert, in dem anderen Buch ist eine einfachere :)
wobei $\alpha = h \nu / m_0 c²$ und $r_0$ der klassische Elektronenradius ist (vgl. \cite{book:knoll}).
Der Energieübertrag auf das Elektron wird berechnet durch
\begin{align}
    E_{e} = E_{\gamma} \left\{1- \frac{1}{1+ E_{\gamma}(1-\cos\theta)/m_0 c^2} \right\}
    \label{eq:ecompton}
\end{align} 


Bei Energien oberhalb der doppelten Ruhemasse des Elektrons (\qty{1.02}{\MeV}) ist die Elektronenpaarbildung möglich.
Diese Wechselwirkung tritt aber realistischerweise nur bei sehr hochenergetischen Gammastrahlen auf.
%Die Dunkelheit ist die Ursache des Lichts.


\subsection{Wahrscheinlichste Photonenenergien}
\begin{itemize} 
    \item[\ce   {^{152}Eu}] - \qty{ 40.1186 } {\keV} \qty{37.7 (5)}{\%} 
    \item[\ce   {^{152}Eu}] - \qty{121.7817 (3)} {\keV} \qty{28.41 (13)}{\%} 
    \item[\ce   {^{152}Eu}] - \qty{344.2785 (12)} {\keV} \qty{26.59 (12)}{\%}
    \item[\ce   {^{137}Cs}] - \qty{661.6553 (30)} {\keV} \qty{85.01 (20)}{\%}
    \item[\ce   {^{133}Ba}] - \qty{30.9731}{\keV} \qty{62.4 (7)}{\%} x ray peak
    \item[\ce   {^{133}Ba}] - \qty{356.0129 (7)}{\keV} \qty{62.05 (19) 	}{\%} breiterer peak
    \item[\ce   {^{125}Sb}] - \qty{427.88}{\keV} \qty{29.6}{\%}
\end{itemize}
\cite{web:lara}

\subsection{Enegieübertrag in einem Detektor \cite{book:gilmore}}
% Das Spektrum eines Monochromatischen Gammastrahlers wie \textbf{Cs} ist in eine 
Alle Signifikanten Interaktionen von Gamma Strahlen mit Materie führen dazu, dass die Energie des Photons 
an ein Elektron weitergegeben wird. 
Dieses Elektron kann in einem Germanium Detektor Elektronenlochpaare erzeugen, wenn diese die notwendige Energie von
$\epsilon = \qty{2.96}{\eV}$ erreicht. 
Die Anzahl an Elektronenlochpaaren ist mit 
\begin{align}
    N = E_e / \epsilon 
\end{align}
proportional zu der Energie $E_e$ des Elektrons.
Basierend auf der Stärke des Halbleiterstroms lassen sich so Rückschlüsse auf die Energie des Elektrons ziehen.
% Gamma Strahlen Photoeffekt -> Voller Energieübertrag
% Compton: Energie basierend auf theta 

\subsection{Form des Detektors}
Der Detektor mit dem wir arbeiten ist eher klein.
Das führt dazu, dass die meisten Compton gestreuten Gammastrahlen meist nur einmal ihre Energie in dem
Germanium Detektor hinterlassen und danach aus dem Detektor heraus reflektieren.
Das Ergebnis ist ein Spektrum in der Form des Compton Kontinuums vgl

\subsection{Fehlerrechnung}
Für die Fehlerrechnung werden alle \textbf{Mittelwerte} von $N$ Messungen folgendermaßen berechnet:

\begin{equation}
    \overline{x} = \frac{1}{N} \cdot \sum_{i=1}^N x_i
    \label{eqn:Mittelwert}
\end{equation}

und alle \textbf{Standardabweichungen zum Mittelwert} mit:

\begin{equation}
    \increment\overline{x} = \sqrt{\frac{1}{N\cdot(N-1)}\cdot\sum_{i=1}^N (x_i-\overline{x})^2}
    \label{eqn:St_Mittelwert}
\end{equation}

Der Fehler für zusammenhängende Messwerte wird dann mit der \textbf{Gaußschen Fehlerfortpflanzung} berechnet:

\begin{equation}
    \increment{f} = \sqrt{ \sum_{i = 1}^{N}  \biggl(\frac{\partial{f}}{\partial{x_i}}\biggr)^2\cdot(\increment{x_i})^2}
    \label{eqn:Gauss}
\end{equation}

Die Fehlerfortpflanzung wird mit Uncertainties in Python \cite{uncertainties} ermittelt.

%---------------------------------------------------------------------------------------------------------------------------------------------------------------%

\section{Durchführung}


%---------------------------------------------------------------------------------------------------------------------------------------------------------------%

\section{Auswertung}


%---------------------------------------------------------------------------------------------------------------------------------------------------------------%

\section{Diskussion}


%---------------------------------------------------------------------------------------------------------------------------------------------------------------%
\newpage
\printbibliography

\end{document}